\chapter{طبقه‌بندی: مفاهیم پایه}
\label{ch:classification-basic}

\section*{اهداف این فصل}
\begin{itemize}
    \item درک مفهوم یادگیری با نظارت و مسئلهٔ طبقه‌بندی
    \item یادگیری نحوهٔ ساخت درخت تصمیم و معیارهای انتخاب ویژگی
    \item آشنایی با هرس درخت تصمیم برای جلوگیری از بیش‌برازش
    \item یادگیری طبقه‌بند بیز ساده
    \item آشنایی مقدماتی با ماتریس درهم‌ریختگی و معیار دقت
\end{itemize}

\section{مقدمه‌ای بر طبقه‌بندی}
\label{sec:classification-intro}

\dmterm{\textbf{طبقه‌بندی}}{Classification} یکی از رایج‌ترین وظایف
\dmterm{یادگیری با نظارت}{Supervised Learning} است. در این نوع یادگیری،
برخلاف خوشه‌بندی که در فصل \ref{ch:clustering} خواهیم دید، مجموعه‌ای از
نمونه‌های \dmterm{\textbf{برچسب‌دار}}{Labeled Data} در اختیار داریم؛ یعنی
برای هر نمونه، علاوه بر ویژگی‌های ورودی، مقدار کلاس هدف نیز مشخص است. هدف
طبقه‌بندی، یادگیری یک تابع $f: X \rightarrow Y$ است که بتواند برای
نمونه‌های جدید و دیده‌نشده، مقدار کلاس $Y$ را از روی ویژگی‌های $X$
پیش‌بینی کند.

مثال‌های رایج طبقه‌بندی عبارت‌اند از: تشخیص اسپم بودن یک ایمیل (بله/خیر)،
پیش‌بینی نکول وام یک متقاضی (پرداخت/نکول)، و تشخیص نوع گل بر اساس اندازهٔ
گلبرگ‌ها (مجموعه دادهٔ معروف \LR{Iris}).

فرآیند کلی ساخت یک مدل طبقه‌بندی دو مرحله دارد:

\begin{enumerate}
    \item \dmterm{\textbf{مرحلهٔ آموزش}}{Training}: مدل با استفاده از
          \dmterm{مجموعهٔ آموزش}{Training Set} که شامل نمونه‌های
          برچسب‌دار است، الگوی رابطهٔ میان ویژگی‌ها و کلاس را یاد می‌گیرد.
    \item \dmterm{\textbf{مرحلهٔ آزمون}}{Testing}: عملکرد مدل روی مجموعه‌ای
          جداگانه به نام \dmterm{مجموعهٔ آزمون}{Test Set} که در آموزش دیده
          نشده، سنجیده می‌شود تا از توانایی
          \dmterm{تعمیم}{Generalization} مدل مطمئن شویم.
\end{enumerate}

\section{درخت تصمیم}
\label{sec:decision-tree}

\dmterm{\textbf{درخت تصمیم}}{Decision Tree} یکی از قابل‌فهم‌ترین و
پرکاربردترین الگوریتم‌های طبقه‌بندی است. ساختار آن شبیه یک فلوچارت است:
هر گره داخلی یک آزمون روی یک ویژگی است، هر شاخه نتیجهٔ آن آزمون را نشان
می‌دهد و هر \dmterm{برگ}{Leaf} یک برچسب کلاس است. مزیت اصلی درخت تصمیم،
تفسیرپذیری بالای آن است؛ می‌توان مسیر از ریشه تا هر برگ را به‌صورت یک
قانون \LR{IF-THEN} ساده خواند.

سؤال اصلی در ساخت درخت تصمیم این است: در هر گره، کدام ویژگی برای تقسیم
داده انتخاب شود؟ پاسخ به این سؤال با معیارهای زیر داده می‌شود.

\subsection{معیارهای انتخاب ویژگی}
\label{subsec:split-criteria}

\paragraph{آنتروپی و بهرهٔ اطلاعاتی}
\dmterm{\textbf{بهرهٔ اطلاعاتی}}{Information Gain} معیاری است که در
الگوریتم \LR{ID3} استفاده می‌شود و بر پایهٔ مفهوم آنتروپی از نظریهٔ
اطلاعات ساخته شده است. برای یک مجموعه داده $D$ با $m$ کلاس، آنتروپی
به‌صورت زیر تعریف می‌شود:
\[
\text{Entropy}(D) = -\sum_{i=1}^{m} p_i \log_2(p_i)
\]
که در آن $p_i$ نسبت نمونه‌های کلاس $i$ در $D$ است. آنتروپی معیاری از
«ناخالصی» یا عدم قطعیت مجموعه داده است: هرچه کلاس‌ها یکنواخت‌تر توزیع
شده باشند، آنتروپی بیشتر است؛ اگر همهٔ نمونه‌ها به یک کلاس تعلق داشته
باشند، آنتروپی برابر صفر است.

پس از تقسیم داده بر اساس ویژگی $A$ به زیرمجموعه‌های $D_1, \ldots, D_v$،
بهرهٔ اطلاعاتی به‌صورت کاهش آنتروپی تعریف می‌شود:
\[
\text{Gain}(D, A) = \text{Entropy}(D) - \sum_{j=1}^{v} \frac{|D_j|}{|D|} \, \text{Entropy}(D_j)
\]
در هر گره، ویژگی‌ای انتخاب می‌شود که بیشترین بهرهٔ اطلاعاتی را دارد.

\paragraph{نسبت بهره}
یک ضعف بهرهٔ اطلاعاتی این است که به‌طور طبیعی به ویژگی‌هایی با تعداد
مقادیر زیاد (مانند شناسهٔ مشتری) گرایش پیدا می‌کند، حتی اگر آن ویژگی از
نظر پیش‌بینی بی‌ارزش باشد. الگوریتم \LR{C4.5} برای رفع این مشکل، معیار
\dmterm{\textbf{نسبت بهره}}{Gain Ratio} را معرفی کرد که بهرهٔ اطلاعاتی را
بر یک عامل نرمال‌ساز به نام
\dmterm{اطلاعات تقسیم}{Split Information} تقسیم می‌کند:
\[
\text{GainRatio}(D, A) = \frac{\text{Gain}(D, A)}{\text{SplitInfo}(D, A)}
\]

\paragraph{شاخص جینی}
\dmterm{\textbf{شاخص جینی}}{Gini Index} معیاری است که در الگوریتم
\LR{CART} استفاده می‌شود و احتمال طبقه‌بندی نادرست یک نمونهٔ تصادفی را
می‌سنجد:
\[
\text{Gini}(D) = 1 - \sum_{i=1}^{m} p_i^2
\]
مشابه آنتروپی، هرچه شاخص جینی کوچک‌تر باشد، مجموعه خالص‌تر است. در هر
گره، ویژگی‌ای انتخاب می‌شود که بیشترین کاهش شاخص جینی را ایجاد کند.

\subsection{هرس درخت تصمیم}
\label{subsec:pruning}

اگر اجازه داده شود درخت تا رسیدن به خلوص کامل هر برگ رشد کند، مدل به
جزئیات و نویز داده‌های آموزش بیش از حد وابسته می‌شود؛ پدیده‌ای که
\dmterm{\textbf{بیش‌برازش}}{Overfitting} نام دارد و باعث افت عملکرد روی
داده‌های جدید می‌شود. \dmterm{\textbf{هرس}}{Pruning} تکنیکی برای مقابله
با این مشکل است و به دو صورت انجام می‌شود:

\begin{itemize}
    \item \dmterm{\textbf{هرس پیشین}}{Pre-pruning}: رشد درخت زودتر متوقف
          می‌شود، مثلاً با تعیین حداکثر عمق درخت یا حداقل تعداد نمونه در
          هر برگ.
    \item \dmterm{\textbf{هرس پسین}}{Post-pruning}: ابتدا اجازه داده
          می‌شود درخت کامل رشد کند، سپس شاخه‌هایی که تأثیر چندانی بر دقت
          روی \dmterm{دادهٔ اعتبارسنجی}{Validation Set} ندارند، حذف
          می‌شوند.
\end{itemize}

در عمل، هرس پسین معمولاً نتایج بهتری می‌دهد، هرچند از نظر محاسباتی
پرهزینه‌تر است.

\section{طبقه‌بند بیز ساده}
\label{sec:naive-bayes}

\dmterm{\textbf{طبقه‌بند بیز ساده}}{Naive Bayes} روشی احتمالاتی برای
طبقه‌بندی است که بر پایهٔ قضیهٔ بیز عمل می‌کند:
\[
P(C \mid X) = \frac{P(X \mid C) \, P(C)}{P(X)}
\]
که در آن $C$ کلاس هدف و $X = (x_1, x_2, \ldots, x_n)$ بردار ویژگی‌های
یک نمونه است. هدف، یافتن کلاسی است که $P(C \mid X)$ را بیشینه می‌کند.

نام «ساده» (\LR{Naive}) به این دلیل است که این روش فرض می‌کند ویژگی‌ها،
مشروط بر کلاس، از یکدیگر \textbf{مستقل} هستند. با این فرض:
\[
P(X \mid C) = \prod_{i=1}^{n} P(x_i \mid C)
\]
با وجود ساده‌بودن (و اغلب نادرست‌بودن) این فرض استقلال در دنیای واقعی،
بیز ساده در عمل، به‌ویژه در طبقه‌بندی متن (مانند تشخیص اسپم)، عملکرد
شگفت‌آور خوبی دارد و از نظر محاسباتی بسیار کارآمد است، چون نیازی به
یادگیری روابط پیچیده میان ویژگی‌ها ندارد.

\paragraph{مسئلهٔ احتمال صفر و هموارسازی لاپلاس}
اگر یک مقدار ویژگی در دادهٔ آموزش هرگز با یک کلاس خاص هم‌رخداد نداشته
باشد، $P(x_i \mid C)$ برابر صفر محاسبه می‌شود و کل حاصل‌ضرب را صفر
می‌کند، حتی اگر سایر ویژگی‌ها به‌شدت بر آن کلاس دلالت کنند. برای رفع این
مشکل، از \dmterm{\textbf{هموارسازی لاپلاس}}{Laplace Smoothing} استفاده
می‌شود که یک مقدار کوچک به همهٔ شمارش‌ها اضافه می‌کند تا هیچ احتمالی
دقیقاً صفر نشود.

\section{ارزیابی مقدماتی}
\label{sec:basic-evaluation}

پس از آموزش یک مدل طبقه‌بندی، باید عملکرد آن را روی دادهٔ آزمون بسنجیم.
ابزار پایه برای این کار
\dmterm{\textbf{ماتریس درهم‌ریختگی}}{Confusion Matrix} است که برای یک
مسئلهٔ دوکلاسه به شکل زیر است:

\begin{center}
\begin{tabular}{@{}lcc@{}}
\toprule
 & \textbf{پیش‌بینی: مثبت} & \textbf{پیش‌بینی: منفی} \\
\midrule
\textbf{واقعی: مثبت} & \LR{True Positive (TP)} & \LR{False Negative (FN)} \\
\textbf{واقعی: منفی} & \LR{False Positive (FP)} & \LR{True Negative (TN)} \\
\bottomrule
\end{tabular}
\end{center}

ساده‌ترین معیار برخاسته از این ماتریس،
\dmterm{\textbf{دقت کلی}}{Accuracy} است:
\[
\text{Accuracy} = \frac{TP + TN}{TP + TN + FP + FN}
\]

هرچند دقت کلی معیاری شهودی است، اما در مسائلی که کلاس‌ها نامتوازن‌اند
(مثلاً تشخیص تقلب که تنها ۱ درصد تراکنش‌ها متقلبانه‌اند) می‌تواند
گمراه‌کننده باشد؛ مدلی که همیشه «غیرمتقلبانه» پیش‌بینی کند، دقت ۹۹ درصد
دارد اما کاملاً بی‌فایده است. معیارهای دقیق‌تر مانند \LR{Precision}،
\LR{Recall} و \LR{F1-Score}، که این مشکل را برطرف می‌کنند، در فصل
\ref{ch:classification-advanced} به‌طور کامل بررسی خواهند شد.

\section*{پیاده‌سازی در پایتون}

پیاده‌سازی درخت تصمیم (\texttt{DecisionTreeClassifier}) با معیارهای
\texttt{entropy} و \texttt{gini}، و همچنین طبقه‌بند \texttt{GaussianNB}،
به‌همراه رسم ماتریس درهم‌ریختگی، در بخش «فصل ۳» پیوست الف
(\ref{app:python}) آمده است.

\section*{خلاصهٔ فصل}

در این فصل با مفهوم یادگیری با نظارت و طبقه‌بندی آشنا شدیم. درخت تصمیم
را به‌عنوان یک مدل قابل‌تفسیر بررسی کردیم و سه معیار اصلی انتخاب ویژگی
(بهرهٔ اطلاعاتی، نسبت بهره و شاخص جینی) و همچنین اهمیت هرس برای جلوگیری
از بیش‌برازش را دیدیم. سپس طبقه‌بند بیز ساده را به‌عنوان یک روش
احتمالاتی و کارآمد معرفی کردیم. در پایان، با ماتریس درهم‌ریختگی و معیار
دقت، اولین ابزار ارزیابی مدل را آموختیم. در فصل بعد، به سراغ
الگوریتم‌های پیشرفته‌تر طبقه‌بندی (نزدیک‌ترین همسایه، ماشین بردار
پشتیبان، روش‌های ترکیبی) و معیارهای ارزیابی دقیق‌تر می‌رویم.

\section*{تمرین‌ها}

\begin{enumerate}
    \item آنتروپی یک مجموعه داده با ۱۰ نمونه که ۶ نمونه از کلاس مثبت و
          ۴ نمونه از کلاس منفی هستند را محاسبه کنید.
    \item چرا معیار بهرهٔ اطلاعاتی به‌تنهایی می‌تواند به ویژگی‌هایی با
          مقادیر بسیار زیاد (مانند شناسهٔ یکتای مشتری) گرایش پیدا کند؟
          نسبت بهره چگونه این مشکل را حل می‌کند؟
    \item تفاوت هرس پیشین و هرس پسین را از نظر هزینهٔ محاسباتی و کیفیت
          نتیجهٔ نهایی مقایسه کنید.
    \item فرض استقلال شرطی در بیز ساده را با یک مثال نقض کنید (یعنی
          مثالی از دو ویژگی که مستقل از هم نیستند) و توضیح دهید چرا با
          وجود این نقض، مدل همچنان می‌تواند عملکرد قابل‌قبولی داشته باشد.
    \item برای یک مسئلهٔ تشخیص بیماری نادر که تنها ۲ درصد جمعیت به آن
          مبتلا هستند، توضیح دهید چرا دقت کلی به‌تنهایی معیار مناسبی
          برای ارزیابی مدل نیست.
\end{enumerate}
